% 5.2 章 · 反问题与正则化
% 这一章的笔记要把"解出来"拆成四件互不相同的事：
%   1) 算子本身丢了什么信息（奇异值衰减、零空间、离散 Picard 条件）；
%   2) 用什么滤波器把剩下的信息取出来（Tikhonov / TSVD / Landweber / TV）；
%   3) 正则参数按什么判据选，以及不同判据之间差多少（oracle / L 曲线 / GCV / 偏差原则）；
%   4) 正则化的代价：偏差与方差此消彼长，两侧各是一种失败方式。
% 所有结论都配一条带实测数字的断言，附表由 tools/build_book.py 生成后 \input 进来。

\numlabchapter{反问题与正则化}{5.2}

\section{这一章要解决什么}

给定 $A\in\mathbb{R}^{m\times n}$ 与带噪数据 $b=Ax^\star+e$，求 $x^\star$。
如果 $A$ 的奇异值衰减得快（或直接秩亏），"最小二乘解" $x_{\mathrm{ls}}=A^{+}b$ 会被噪声吃光。
本章不把 Tikhonov、TSVD、Landweber、TV 当成四张要背的公式，而是把下面五句话算清楚：

\begin{enumerate}
  \item \textbf{不做正则化会错多少？} 误差被放大 $\lVert e\rVert/\sigma_n$；$\kappa=10^6$、
        相对噪声 $10^{-3}$ 时最坏放大实测 $1028$ 倍、均方根 $182$ 倍（定理~\ref{thm:amplify}）。
  \item \textbf{什么时候"截断"这件事是有意义的？} 离散 Picard 条件：$\lvert u_i^{\mathsf T}b\rvert$
        必须比 $\sigma_i$ 衰减得更快。$\mathrm{decay}=2$ 的无噪算例上斜率实测 $2.000000$，
        噪声 $10^{-3}$ 时整段斜率被压到 $0.564$，而\emph{信号段}仍是 $1.997$
        （定理~\ref{thm:picard}）。
  \item \textbf{"最优 $\lambda$"是按哪个判据最优？} $\kappa=10^6$、噪声 $10^{-3}$ 时，
        oracle $1.585\times10^{-2}$、L 曲线 $7.36\times10^{-4}$、GCV $7.36\times10^{-2}$、
        偏差原则 $1.34\times10^{-2}$——误差放大分别是 $1.0$、$21.5$、$3.79$、$1.16$ 倍，
        完全不正则化是 $12863$ 倍（定理~\ref{thm:discrepancy}）。
  \item \textbf{三种滤波器可以公平比较吗？} 只能在同一条\emph{有效自由度}
        $\sum_i f_i$ 轴上比：自由度对齐到 $8$ 时 TSVD $3.50\times10^{-3}$、
        Tikhonov $2.83\times10^{-3}$、Landweber $9.24\times10^{-2}$
        （定理~\ref{thm:landweber}）。
  \item \textbf{正则化一定更准吗？} 不是。$\lambda\to0$ 时噪声接管（实测误差 $1.53$），
        $\lambda\to\infty$ 时解被压成 $0$（实测误差 $1.27$，回到 $\lVert x^\star\rVert$），
        中间才有最小值 $0.0199$（定理~\ref{thm:biasvar}）。
\end{enumerate}

记号：$A=U\operatorname{diag}(\sigma_1,\dots,\sigma_r)V^{\mathsf T}$ 是 1.3 章的
thin SVD（$U\in\mathbb{R}^{m\times r}$、$V\in\mathbb{R}^{n\times r}$、$r=\min(m,n)$，
$\sigma_1\ge\dots\ge\sigma_r>0$），$\kappa=\sigma_1/\sigma_r$；
$e$ 是噪声、$\delta=\lVert e\rVert$；$\alpha_i=u_i^{\mathsf T}b$；
$f_i$ 是滤波因子；$\operatorname{dof}(\lambda)=\sum_i f_i$ 是有效自由度；
$u=\tfrac12\varepsilon_{\text{mach}}\approx1.11\times10^{-16}$。

\section{病态性与噪声放大}

\begin{theorem}[最小二乘解的噪声放大与等号条件]
\label{thm:amplify}
设 $A$ 满列秩，$x_{\mathrm{ls}}=A^{+}b$，则
\[
  \lVert x_{\mathrm{ls}}-x^\star\rVert_2=\lVert A^{+}e\rVert_2
  \le \frac{\lVert e\rVert_2}{\sigma_r},
\]
且等号成立当且仅当 $e$ 落在 $u_r$ 张成的一维子空间上。
若 $e$ 的各分量独立同分布、方差为 $\delta^2/m$，则
\[
  \mathbb{E}\lVert A^{+}e\rVert_2^2=\frac{\delta^2}{m}\sum_{i=1}^{r}\frac{1}{\sigma_i^2},
  \qquad
  \frac{\text{最坏放大}}{\text{均方根放大}}
  =\frac{\delta/\sigma_r}{\delta\sqrt{\sum_i\sigma_i^{-2}}/\sqrt{m}}
  =\frac{\sqrt{m}}{\sigma_r\sqrt{\sum_i\sigma_i^{-2}}}\ \le\ \sqrt{m},
\]
最后的不等号在 $\sigma_r^{-2}$ 主导求和时取到等号（几何衰减谱、$\kappa$ 大时就是这种情形）。
\end{theorem}

\begin{proof}
把 $e$ 在 $\{u_i\}$ 下展开：$e=\sum_i (u_i^{\mathsf T}e)u_i+w$，其中 $w\perp\operatorname{range}(U)$。
由 $A^{+}=\sum_i \sigma_i^{-1}v_iu_i^{\mathsf T}$ 得
$A^{+}e=\sum_i \sigma_i^{-1}(u_i^{\mathsf T}e)v_i$，而 $A^{+}w=0$。
于是 $\lVert A^{+}e\rVert_2^2=\sum_i \sigma_i^{-2}(u_i^{\mathsf T}e)^2$。
把每一项放缩为 $\sigma_i^{-2}(u_i^{\mathsf T}e)^2\le\sigma_r^{-2}(u_i^{\mathsf T}e)^2$ 并求和，
再用 $\sum_i(u_i^{\mathsf T}e)^2\le\lVert e\rVert^2$（$w$ 不贡献）即得最坏界；
取 $e=\lVert e\rVert u_r$ 时 $\sum_i\sigma_i^{-2}(u_i^{\mathsf T}e)^2=\lVert e\rVert^2\sigma_r^{-2}$，等号达成。
期望式由 $\mathbb{E}(u_i^{\mathsf T}e)^2=\delta^2/m$（正交不变性）直接得到。
\end{proof}

实测（$\kappa=10^6$、$\delta/\lVert b_{\text{exact}}\rVert=10^{-3}$、$32\times20$）：
噪声对齐 $u_r$ 时相对误差 $900.363172$，解析值 $\delta/(\sigma_r\lVert x^\star\rVert)$ 也是
$900.363172$（相对差 $7.9\times10^{-12}$）；最坏放大 $1028.4$、均方根 $181.8$，
两者之比 $5.66=\sqrt{m}$ 的 $1.0$ 倍——因为这时 $\sum_i\sigma_i^{-2}$ 由 $\sigma_r^{-2}$ 主导。

\begin{theorem}[离散 Picard 条件与截断的必要性]
\label{thm:picard}
若存在常数 $C>0$、指数 $p>1$ 使 $\lvert u_i^{\mathsf T}b_{\text{exact}}\rvert\le C\sigma_i^{p}$，
则真解系数 $\lvert v_i^{\mathsf T}x^\star\rvert=\lvert u_i^{\mathsf T}b_{\text{exact}}\rvert/\sigma_i
\le C\sigma_i^{p-1}$ 随 $i$ 递减；反之若 $p<1$，系数随 $i$ \emph{增长}，
此时 $x^\star$ 的范数中来自小奇异值方向的分量占主导。
\end{theorem}

\begin{proof}
由 $b_{\text{exact}}=Ax^\star=U\operatorname{diag}(\sigma)V^{\mathsf T}x^\star$ 得
$u_i^{\mathsf T}b_{\text{exact}}=\sigma_i(v_i^{\mathsf T}x^\star)$，即
$v_i^{\mathsf T}x^\star=\sigma_i^{-1}(u_i^{\mathsf T}b_{\text{exact}})$。
把假设代入即得 $\lvert v_i^{\mathsf T}x^\star\rvert\le C\sigma_i^{p-1}$；
$p<1$ 时 $\sigma_i^{p-1}=\sigma_i^{-(1-p)}$ 随 $i$ 增大而增大。
在双对数坐标下，"系数对 $\sigma_i$ 的斜率 $p$"与"真解系数对 $\sigma_i$ 的斜率 $p-1$"
是同一个事实的两种写法：$p>1$ 等价于真解系数递减。
\end{proof}

实测：造题时取 $\alpha_i=\sigma_i^{2}$，无噪时整段斜率 $2.000000$、比值斜率 $1.000000$、
相关系数 $1.000000$；噪声 $10^{-3}$ 时整段斜率 $0.564$、信号段（前 $5$ 个模态）$1.997$，
$20$ 个模态里 $15$ 个的 $\lvert u_i^{\mathsf T}b\rvert$ 已经落到噪声之下。
把造题指数改成 $0.5$（Picard 不成立）后，真误差最小点移到 $\lambda$ 网格左端，
"最优正则化解"与完全不正则化误差相同（放大 $1.000$ 倍）。

\section{Tikhonov 与谱滤波}

\begin{theorem}[Tikhonov 解的滤波因子形式]
\label{thm:tikhonov}
$A$ 的 thin SVD 如上。对任意 $\lambda>0$，
\[
  x_\lambda=\arg\min_x\ \lVert Ax-b\rVert_2^2+\lambda^2\lVert x\rVert_2^2
  =\sum_{i=1}^{r} f_i(\lambda)\,\frac{u_i^{\mathsf T}b}{\sigma_i}\,v_i,
  \qquad f_i(\lambda)=\frac{\sigma_i^{2}}{\sigma_i^{2}+\lambda^{2}} .
\]
并且：(i) $0<f_i<1$，$f_i$ 对 $\lambda$ 严格递减；(ii) $f_i=1/2$ 当且仅当 $\sigma_i=\lambda$；
(iii) $\lVert x_\lambda\rVert$ 对 $\lambda$ 不增，$\lVert Ax_\lambda-b\rVert$ 对 $\lambda$ 不减，
$\operatorname{dof}(\lambda)=\sum_if_i$ 对 $\lambda$ 不增且取值于 $(0,r)$。
\end{theorem}

\begin{proof}
目标函数对 $x$ 求梯度并令其为零：$(A^{\mathsf T}A+\lambda^2I)x=A^{\mathsf T}b$。
代入 $A=U\Sigma V^{\mathsf T}$ 得 $V(\Sigma^2+\lambda^2I)V^{\mathsf T}x=V\Sigma U^{\mathsf T}b$，
即 $x_\lambda=V(\Sigma^2+\lambda^2I)^{-1}\Sigma U^{\mathsf T}b
=\sum_i\frac{\sigma_i}{\sigma_i^2+\lambda^2}(u_i^{\mathsf T}b)v_i$，
把 $\sigma_i/(\sigma_i^2+\lambda^2)$ 写成 $\frac{1}{\sigma_i}\cdot\frac{\sigma_i^2}{\sigma_i^2+\lambda^2}$
即得滤波因子形式。(i) 与 (ii) 是对 $\lambda$ 求导与令 $f=1/2$ 的直接结果。
(iii)：$\lVert x_\lambda\rVert^2=\sum_i f_i^2\alpha_i^2/\sigma_i^2$，
逐项 $\mathrm{d}f_i/\mathrm{d}\lambda<0$ 说明它不增；
残差 $\lVert Ax_\lambda-b\rVert^2=\sum_i(1-f_i)^2\alpha_i^2+\lVert P_\perp b\rVert^2$，
$1-f_i=\lambda^2/(\sigma_i^2+\lambda^2)$ 对 $\lambda$ 递增；
自由度是单调递减函数之和。
\end{proof}

实测：滤波因子与闭式的最大偏差 $0.0$，$f(\sigma=\lambda)=0.500000000000$；
$\lambda=0$ 时与 1.3 章 \texttt{least\_squares\_svd} 的解相对差 $2.8\times10^{-11}$
（$\kappa=10^6$ 下不得不付出的抵消误差）；
谱公式与用 LU 解正规方程 $(A^{\mathsf T}A+\lambda^2I)x=A^{\mathsf T}b$ 的结果，
在 $\lambda=10^{-2}$ 时相对差 $2.0\times10^{-12}$，而 $\lambda=10^{-6}$ 时退化到
$3.2\times10^{-6}$（$\kappa(A)^2=10^{12}$，正规方程路线自己丢了 $6$ 位）。

\begin{theorem}[偏差-方差分解与方差的两段幂律]
\label{thm:biasvar}
设 $b=b_{\text{exact}}+e$，$\mathbb{E}e=0$、$\mathbb{E}ee^{\mathsf T}=(\delta^2/m)I_m$，
$x_\lambda$ 为 Tikhonov 解，则
\[
  \mathbb{E}\lVert x_\lambda-x^\star\rVert_2^2
  =\underbrace{\sum_{i=1}^{r}(f_i-1)^2\frac{\alpha_i^2}{\sigma_i^2}}_{\text{偏差}^2}
  +\underbrace{\frac{\delta^2}{m}\sum_{i=1}^{r}\frac{f_i^2}{\sigma_i^2}}_{\text{方差}} .
\]
当 $\lambda\ge10\sigma_1$ 时方差满足 $\text{方差}=\Theta(\lambda^{-4})$；
当 $\lambda\le\sigma_r$ 时方差趋于常数 $\frac{\delta^2}{m}\sum_i\sigma_i^{-2}$。
\end{theorem}

\begin{proof}
$b$ 只通过 $\alpha_i=u_i^{\mathsf T}b$ 进入解，而 $\mathbb{E}\alpha_i=u_i^{\mathsf T}b_{\text{exact}}$、
$\operatorname{Var}(\alpha_i)=\delta^2/m$。于是
$\mathbb{E}x_\lambda=\sum_i f_i\frac{\mathbb{E}\alpha_i}{\sigma_i}v_i$，偏差部分为
$\sum_i(f_i-1)\frac{u_i^{\mathsf T}b_{\text{exact}}}{\sigma_i}v_i$，取范数平方即得第一项。
方差部分：各 $\alpha_i$ 独立，$\operatorname{Var}(x_\lambda)=\sum_i\frac{f_i^2}{\sigma_i^2}\operatorname{Var}(\alpha_i)v_iv_i^{\mathsf T}$，
取迹得 $\frac{\delta^2}{m}\sum_if_i^2/\sigma_i^2$。渐近式：$\lambda\gg\sigma_1$ 时
$f_i\approx\sigma_i^2/\lambda^2$，故 $f_i/\sigma_i\approx\sigma_i/\lambda^2$，
方差 $\approx\frac{\delta^2}{m\lambda^4}\sum_i\sigma_i^2=\Theta(\lambda^{-4})$；
$\lambda\ll\sigma_r$ 时 $f_i\approx1$，方差 $\to\frac{\delta^2}{m}\sum_i\sigma_i^{-2}$。
\end{proof}

实测（$\kappa=10^4$、噪声 $10^{-3}$）：解析与 $64$ 个样本的蒙特卡洛总误差最大相对差
$0.099$、方差 $0.118$（$1/\sqrt{64}=0.125$ 量级），偏差$^2+$方差$=$总误差逐位成立；
大方差端的对数-对数斜率 $-3.9993$（理论 $-4$）；$\lambda$ 两端误差 $1.53$ 与 $1.27$，
中间最小值 $0.0199$。注意蒙特卡洛的\emph{偏差}在小 $\lambda$ 端被样本均值污染：
解析偏差$^2=6.4\times10^{-9}$、实测 $4.35\times10^{-2}$（相差 $106.5$ 倍 ${\rm mse}_{\min}$），
所以偏差必须用无噪的 $b_{\text{exact}}$ 算。

\section{L 曲线、GCV 与偏差原则}

\begin{theorem}[残差对 $\lambda$ 的单调性与偏差原则交点的唯一性]
\label{thm:discrepancy}
记 $\rho(\lambda)=\lVert Ax_\lambda-b\rVert_2$。则 $\rho$ 在 $(0,\infty)$ 上连续、
严格递增，且
\[
  \rho(0+)=\lVert P_\perp b\rVert_2,\qquad \rho(\infty)=\lVert b\rVert_2 .
\]
因此若 $\lVert P_\perp b\rVert_2<\tau\delta<\lVert b\rVert_2$（$\tau\ge1$），
方程 $\rho(\lambda)=\tau\delta$ 在 $(0,\infty)$ 上有且只有一个解。
\end{theorem}

\begin{proof}
由定理~\ref{thm:tikhonov} 的推导，
$\rho(\lambda)^2=\lVert P_\perp b\rVert^2+\sum_i\bigl(\frac{\lambda^2}{\sigma_i^2+\lambda^2}\bigr)^2\alpha_i^2$。
每一项 $\bigl(\lambda^2/(\sigma_i^2+\lambda^2)\bigr)^2\alpha_i^2$ 关于 $\lambda$ 严格递增
（$\lambda^2/(\sigma_i^2+\lambda^2)=1-\sigma_i^2/(\sigma_i^2+\lambda^2)$，后一项严格递减），
有限个严格递增函数之和仍严格递增，故 $\rho$ 严格递增、连续。
$\lambda\to0$ 时每项 $\to0$，$\lambda\to\infty$ 时每项 $\to\alpha_i^2$，
于是 $\rho(\infty)^2=\lVert P_\perp b\rVert^2+\sum_i\alpha_i^2=\lVert b\rVert^2$。
严格单调连续函数的像是一个区间，$(\rho(0+),\rho(\infty))$ 中的每个值恰好取一次。
\end{proof}

实测：$\lambda_{\text{disc}}=1.34\times10^{-2}$ 处
$\rho/\delta=0.985457$（对数插值带来的 $1.5\%$ 偏差，交点落在两个网格点之间）。
注意 $m=32>r=20$ 时 $\lVert P_\perp b\rVert\approx\sqrt{1-r/m}\,\delta=0.61\,\delta$
（$P_\perp b=P_\perp e$），它**低于**目标 $\tau\delta$，所以交点存在；
但当 $\delta$ 小到与拟合/舍入地板同量级时，$\rho$ 的地板已经高于 $\tau\delta$，
方程在数值上无解——帧里这时把 $\lambda_{\text{disc}}$ 写成 \texttt{null} 并在 \texttt{strings} 里说明。

\begin{theorem}[GCV 函数的分母与它为何需要 $m>r$]
\label{thm:gcv}
\[
  G(\lambda)=\frac{\lVert Ax_\lambda-b\rVert_2^2}{\bigl(m-\operatorname{dof}(\lambda)\bigr)^2}
\]
在 $\lambda\to0$ 时趋于 $\lVert P_\perp b\rVert^2/(m-r)^2$（$m>r$ 时有限）；
若 $m=r$ 则分母趋于 $0$、$G$ 发散，$\lambda=0$ 成为一个"无穷小的分母"造出来的伪极小点。
\end{theorem}

\begin{proof}
由定理~\ref{thm:tikhonov}，$\operatorname{dof}(\lambda)=\sum_if_i\to r$（$\lambda\to0$），
故分母 $\to(m-r)^2$；分子 $\to\lVert P_\perp b\rVert^2$（同上）。
$m=r$ 时分母 $\to0$ 而分子 $\to0$（若 $b\in\operatorname{range}(A)$），比值的行为由
两者的阶数决定，实现里必须给分母一个地板，并且**不**把这一段当成有限的最小值。
\end{proof}

实测：$G$ 与 1.3 章 \texttt{cross\_validation\_curve} 的最大相对差 $3.5\times10^{-12}$，
最小值位置一致；$\lambda_{\mathrm{GCV}}$ 对五个噪声档位（$10^{-2}\dots10^{-6}$）的
幂律指数 $0.4667$，oracle 是 $0.4800$——两者都接近 $1/2$，GCV 的滞后量就是它们之差。

\section{TSVD、Landweber 与早停}

\begin{theorem}[Landweber 迭代的谱形式与收敛充要条件]
\label{thm:landweber}
迭代 $x_{k+1}=x_k+\omega A^{\mathsf T}(b-Ax_k)$，$x_0=0$。则
\[
  x_k=\sum_{i=1}^{r}\bigl(1-(1-\omega\sigma_i^2)^k\bigr)\frac{u_i^{\mathsf T}b}{\sigma_i}v_i ,
\]
并且 $\{x_k\}$ 对任意初值收敛的充要条件是
\[
  0<\omega<\frac{2}{\sigma_1^2}.
\]
此时滤波因子 $f_i^{(k)}=1-(1-\omega\sigma_i^2)^k$ 对 $k$ 单调不减、取值于 $[0,1)$；
$\omega=1/\sigma_1^2$ 时最慢模态一步压到 $0$。
\end{theorem}

\begin{proof}
误差递推：$x_{k+1}-x^\star=(I-\omega A^{\mathsf T}A)(x_k-x^\star)+\omega A^{\mathsf T}e$
的齐次部分在 SVD 基下解耦，$V^{\mathsf T}(x_{k+1}-x_k)$ 的第 $i$ 个分量乘
$(1-\omega\sigma_i^2)$，归纳得滤波因子形式。
收敛性由 $\max_i\lvert1-\omega\sigma_i^2\rvert<1$ 给出，即
$-1<1-\omega\sigma_1^2$（最坏的是最大奇异值）与 $1-\omega\sigma_r^2<1$（$\omega>0$），
合并即 $0<\omega<2/\sigma_1^2$。
单调性：$\omega\sigma_i^2\in(0,2)$ 时 $\lvert1-\omega\sigma_i^2\rvert<1$，故
$(1-\omega\sigma_i^2)^k$ 的绝对值随 $k$ 递减且符号固定，$f_i^{(k)}$ 单调升向 $1$。
\end{proof}

实测：谱形式与真正的矩阵迭代（$A$ 与 $A^{\mathsf T}$ 只作为算子给出）在三组
$(\omega_{\text{rel}},k)$ 上相对差 $2.6\times10^{-13}$ 以内；
$\omega_{\text{rel}}=2.2$ 时矩阵迭代 \texttt{stop='diverged'}，滤波因子最大绝对值 $6.9\times10^{15}$。
二分实测发散阈值：矩阵迭代（$400$ 步预算）$2.0715$，谱判据
$2.0000000000009$——前者偏高 $3.6\%$，因为 $\lvert1-\omega\sigma_1^2\rvert$ 刚过 $1$ 时
增长因子只有 $1+\epsilon$，$400$ 步涨不出可观范数。**两条阈值都要报**：只有谱判据那条
才等于 $2/\sigma_1^2$。

\medskip

TSVD 的滤波因子是阶跃 $f_i=\mathbf{1}\{i\le k\}$，自由度恰好是 $k$；
Tikhonov 的过渡带宽度约为一个 $\lambda$ 的十倍程；Landweber 的过渡带由 $k$ 控制。
把三条误差曲线画在同一条自由度轴上（$n=20$、$\kappa=10^6$、噪声 $10^{-4}$、
对齐到 $k=8$）：TSVD $3.50\times10^{-3}$、Tikhonov $2.83\times10^{-3}$、
Landweber $9.24\times10^{-2}$。**形状不同不等于精度不同**，而"哪个更准"必须实测。

\section{全变差}

\begin{theorem}[两点信号的 TV 去噪闭式解]
\label{thm:tv}
设 $b\in\mathbb{R}^2$，$m=\tfrac12(b_1+b_2)$、$d=\tfrac12(b_2-b_1)$。则
\[
  \operatorname*{arg\,min}_{x\in\mathbb{R}^2}\ \tfrac12\lVert x-b\rVert_2^2+\lambda\lvert x_2-x_1\rvert
  =\bigl(m-\operatorname{soft}(d,\lambda),\;m+\operatorname{soft}(d,\lambda)\bigr),
  \qquad \operatorname{soft}(d,\lambda)=\operatorname{sign}(d)\max(\lvert d\rvert-\lambda,0).
\]
特别地 $\lambda=0$ 时 $x=b$；$\lambda\ge\lvert d\rvert$ 时 $x_1=x_2=m$。
\end{theorem}

\begin{proof}
把变量换成 $(m,d)$ 坐标：$x_1=m-d$、$x_2=m+d$，则
$\tfrac12\lVert x-b\rVert^2=\tfrac12\bigl((m-b_1+d)^2+(m-b_2-d)^2\bigr)$，
对 $m$ 求导得 $m=\tfrac12(b_1+b_2)$（$m$ 与 $d$ 解耦）。
剩下关于 $d$ 的一维问题 $\min_d (d-d_0)^2+2\lambda\lvert d\rvert$（$d_0=\tfrac12(b_2-b_1)$），
其最优性条件是 $0\in 2(d-d_0)+2\lambda\partial\lvert d\rvert$：
$d>0$ 时 $d=d_0-\lambda$，$d<0$ 时 $d=d_0+\lambda$，$d=0$ 时要求 $\lvert d_0\rvert\le\lambda$。
合并即软阈值公式。
\end{proof}

实测：$\lambda=0,0.05,0.3,0.75,2$ 五档上与闭式解的最大逐点偏差 $1.1\times10^{-16}$；
$\lambda\ge\lvert d\rvert$ 时两点差 $1.1\times10^{-16}$。
TV 在 SVD 基下**不是**对角算子：用 $b=u_i$ 探针测出的"有效因子"换一个探针
（$b=(u_1+u_n)/\sqrt2$）就改变 $0.959$，非对角泄漏 $0.144$；
$\lambda=0$ 的正对照是探针读数与 $1$ 的偏差 $5.4\times10^{-14}$。

\section{图像与投影算子}

\begin{theorem}[可分离算子的奇异值]
\label{thm:separable}
设二维算子 $\mathcal{A}(X)=BX B^{\mathsf T}$，$B\in\mathbb{R}^{n\times n}$ 的 SVD 为
$B=U\operatorname{diag}(s)V^{\mathsf T}$。则
\[
  \mathcal{A}=(U\otimes U)\operatorname{diag}(s_is_j)(V\otimes V)^{\mathsf T},
\]
即二维算子的奇异值是全部乘积 $s_is_j$（$i,j=1,\dots,n$），
并且 $\kappa_{2D}=\kappa_{1D}^2$、$\operatorname{rank}(\mathcal{A})=\operatorname{rank}(B)^2$。
\end{theorem}

\begin{proof}
把 $X$ 按列拉直，$\operatorname{vec}(BX B^{\mathsf T})=(B\otimes B)\operatorname{vec}(X)$。
Kronecker 积的 SVD 满足 $(B\otimes B)=(U\otimes U)(\operatorname{diag}(s)\otimes\operatorname{diag}(s))(V\otimes V)^{\mathsf T}$，
而 $\operatorname{diag}(s)\otimes\operatorname{diag}(s)=\operatorname{diag}(s_is_j)$
（按字典序排列）。于是奇异值集合是 $\{s_is_j\}$，
最大与最小之比为 $(s_1/s_n)^2$，秩为非零乘积的个数即 $\operatorname{rank}(B)^2$。
\end{proof}

实测：$n=24$、模糊宽度 $\sigma=2$ 时 $\kappa_{1D}=5911$、$\kappa_{2D}=3.494\times10^7=5911^2$；
二维 Tikhonov 解满足正规方程 $(B^{\mathsf T}B)X(B^{\mathsf T}B)+\lambda^2X=B^{\mathsf T}YB$
（相对残差 $5.1\times10^{-15}$）。

\begin{theorem}[投影算子的秩上界与零空间误差地板]
\label{thm:rank}
设 $A\in\mathbb{R}^{m\times n}$，则 $\operatorname{rank}(A)\le\min(m,n)$，
零空间维数 $\dim\ker A\ge n-\operatorname{rank}(A)$。
若 $x^\star$ 在 $\ker A$ 上的投影非零，则对任意 $\lambda$ 与任意噪声水平
\[
  \lVert x_\lambda-x^\star\rVert_2\ \ge\ \lVert P_{\ker A}x^\star\rVert_2
  -\ \lVert x_\lambda-P_{\operatorname{row}}x_\lambda\rVert_2 ,
\]
特别地当正则化解落在行空间内时误差不可能低于 $\lVert P_{\ker A}x^\star\rVert$。
\end{theorem}

\begin{proof}
秩上界是线性代数基本事实（$A$ 的像由 $m$ 个坐标与 $n$ 个列向量共同限制）。
对任意 $x_\lambda=x_\lambda^{\parallel}+x_\lambda^{\perp}$（行空间 $+$ 零空间分解），
$x^\star=x^{\parallel}_\star+x^\perp_\star$，则
$\lVert x_\lambda-x^\star\rVert\ge\lVert x_\lambda^\perp-x^\perp_\star\rVert
\ge\lVert x^\perp_\star\rVert-\lVert x^\perp_\lambda\rVert$（三角不等式）。
Tikhonov/TSVD 的解都在行空间内（$V$ 的列张成行空间），故 $x^\perp_\lambda=0$，
误差下界就是 $\lVert P_{\ker A}x^\star\rVert$。
\end{proof}

实测：$10\times10$ 幻影（$n^2=100$ 个未知量），$8$ 条角度 $\times$ $10$ 条射线 $=80$ 行，
秩 $80$、零空间 $20$，噪声拖到 $10^{-8}$ 时重建误差仍为 $0.158$；
$14$ 条角度（$140\times100$）时满列秩、$\kappa=17.8$，噪声 $10^{-8}$ 时误差 $4.3\times10^{-8}$。
**"CT 一定病态"在这个尺寸上不成立**：区分"条件数大"与"秩亏"是两件事。

\section{四个可视化怎么读}
\label{sec:read}

\begin{enumerate}
  \item \textbf{L 曲线交互}（\texttt{lcurve\_tradeoff}）：横轴 $\log_{10}\lVert x_\lambda\rVert$、
        纵轴 $\log_{10}\rho(\lambda)$。拖动 $\lambda$ 时同一张图上标出四个判据
        （oracle、L 曲线拐点、GCV、偏差原则）各自挑出的点；副轴是真误差曲线。
        三个拐点判据的 $\lambda$ 会差 $33\sim681$ 倍，这一点必须在图上看得见。
        把 Picard 档位从 \texttt{satisfied} 拖到 \texttt{violated}，可以看到正则化收益掉到 $1.0$ 倍。
  \item \textbf{滤波因子}（\texttt{filter\_shapes}）：同一条 $\sigma$ 轴上画 TSVD 的阶跃、
        Tikhonov 的平滑过渡、Landweber 的渐进过渡；TV 只能给探针读数。
        副轴是三条误差-自由度曲线与"同自由度下谁更准"的匹配点。
        $\omega$ 拖过 $2$ 之后 Landweber 当场发散，图上那一段是夹住并带掩码的占位值。
  \item \textbf{图像反卷积 / 简化 CT}（\texttt{image\_deconvolution}）：四类方法重建 $24\times24$ 幻影；
        拖噪声水平看解什么时候被噪声吞掉，拖迭代步看半收敛。
        切到 CT 后 $8$ 条角度会看到"误差地板由零空间决定"。
  \item \textbf{偏差-方差}（\texttt{bias\_variance}）：偏差$^2$、方差、总误差三条曲线，
        解析与蒙特卡洛并列；$\lambda\to0$ 与 $\lambda\to\infty$ 两侧的失败数字都标出来。
\end{enumerate}

% 本文件由 tools/build_book.py 从 dist/reports/assertions.json 生成，请勿手改。
\begin{longtable}{@{}p{0.06\textwidth}p{0.32\textwidth}p{0.22\textwidth}p{0.16\textwidth}@{}}
\caption{章节 5.2 数值断言（生成时间 2026-10-19T17:44:22-0400）}\\
\toprule
类别 & 断言 & 理论值 & 实测值 \\
\midrule
\endfirsthead
\toprule
类别 & 断言 & 理论值 & 实测值 \\
\midrule
\endhead
\bottomrule
\endlastfoot

B & lcurve\_\allowbreak{}tradeoff/log10\_\allowbreak{}lam 全有限 & 0 个非有限值 & 0 \checkmark \\
B & lcurve\_\allowbreak{}tradeoff/log10\_\allowbreak{}x\_\allowbreak{}norm 全有限 & 0 个非有限值 & 0 \checkmark \\
B & lcurve\_\allowbreak{}tradeoff/log10\_\allowbreak{}residual 全有限 & 0 个非有限值 & 0 \checkmark \\
B & lcurve\_\allowbreak{}tradeoff/error\_\allowbreak{}curve 全有限 & 0 个非有限值 & 0 \checkmark \\
B & lcurve\_\allowbreak{}tradeoff/curvature 全有限 & 0 个非有限值 & 0 \checkmark \\
B & lcurve\_\allowbreak{}tradeoff/dof 全有限 & 0 个非有限值 & 0 \checkmark \\
B & lcurve\_\allowbreak{}tradeoff/gcv 全有限 & 0 个非有限值 & 0 \checkmark \\
B & lcurve\_\allowbreak{}tradeoff/criterion\_\allowbreak{}log10\_\allowbreak{}lam 全有限 & 0 个非有限值 & 0 \checkmark \\
B & lcurve\_\allowbreak{}tradeoff/criterion\_\allowbreak{}error 全有限 & 0 个非有限值 & 0 \checkmark \\
B & lcurve\_\allowbreak{}tradeoff/criterion\_\allowbreak{}valid 全有限 & 0 个非有限值 & 0 \checkmark \\
B & lcurve\_\allowbreak{}tradeoff/corner\_\allowbreak{}log10\_\allowbreak{}lam 全有限 & 0 个非有限值 & 0 \checkmark \\
B & lcurve\_\allowbreak{}tradeoff/picard\_\allowbreak{}sigma 全有限 & 0 个非有限值 & 0 \checkmark \\
B & lcurve\_\allowbreak{}tradeoff/picard\_\allowbreak{}coeff 全有限 & 0 个非有限值 & 0 \checkmark \\
B & lcurve\_\allowbreak{}tradeoff/picard\_\allowbreak{}ratio 全有限 & 0 个非有限值 & 0 \checkmark \\
B & lcurve\_\allowbreak{}tradeoff/picard\_\allowbreak{}below\_\allowbreak{}noise 全有限 & 0 个非有限值 & 0 \checkmark \\
B & lcurve\_\allowbreak{}tradeoff/solution\_\allowbreak{}true 全有限 & 0 个非有限值 & 0 \checkmark \\
B & lcurve\_\allowbreak{}tradeoff/solution\_\allowbreak{}lam 全有限 & 0 个非有限值 & 0 \checkmark \\
B & filter\_\allowbreak{}shapes/index 全有限 & 0 个非有限值 & 0 \checkmark \\
B & filter\_\allowbreak{}shapes/log10\_\allowbreak{}sigma 全有限 & 0 个非有限值 & 0 \checkmark \\
B & filter\_\allowbreak{}shapes/filter\_\allowbreak{}tsvd 全有限 & 0 个非有限值 & 0 \checkmark \\
B & filter\_\allowbreak{}shapes/filter\_\allowbreak{}tikhonov 全有限 & 0 个非有限值 & 0 \checkmark \\
B & filter\_\allowbreak{}shapes/filter\_\allowbreak{}landweber 全有限 & 0 个非有限值 & 0 \checkmark \\
B & filter\_\allowbreak{}shapes/filter\_\allowbreak{}landweber\_\allowbreak{}valid 全有限 & 0 个非有限值 & 0 \checkmark \\
B & filter\_\allowbreak{}shapes/filter\_\allowbreak{}tv\_\allowbreak{}probe 全有限 & 0 个非有限值 & 0 \checkmark \\
B & filter\_\allowbreak{}shapes/filter\_\allowbreak{}tv\_\allowbreak{}probe\_\allowbreak{}mixed 全有限 & 0 个非有限值 & 0 \checkmark \\
B & filter\_\allowbreak{}shapes/dof\_\allowbreak{}tsvd\_\allowbreak{}curve 全有限 & 0 个非有限值 & 0 \checkmark \\
B & filter\_\allowbreak{}shapes/error\_\allowbreak{}tsvd\_\allowbreak{}curve 全有限 & 0 个非有限值 & 0 \checkmark \\
B & filter\_\allowbreak{}shapes/dof\_\allowbreak{}tikhonov\_\allowbreak{}curve 全有限 & 0 个非有限值 & 0 \checkmark \\
B & filter\_\allowbreak{}shapes/error\_\allowbreak{}tikhonov\_\allowbreak{}curve 全有限 & 0 个非有限值 & 0 \checkmark \\
B & filter\_\allowbreak{}shapes/log10\_\allowbreak{}lam\_\allowbreak{}curve 全有限 & 0 个非有限值 & 0 \checkmark \\
B & filter\_\allowbreak{}shapes/dof\_\allowbreak{}landweber\_\allowbreak{}curve 全有限 & 0 个非有限值 & 0 \checkmark \\
B & filter\_\allowbreak{}shapes/error\_\allowbreak{}landweber\_\allowbreak{}curve 全有限 & 0 个非有限值 & 0 \checkmark \\
B & filter\_\allowbreak{}shapes/error\_\allowbreak{}matched 全有限 & 0 个非有限值 & 0 \checkmark \\
B & image\_\allowbreak{}deconvolution/image\_\allowbreak{}true 全有限 & 0 个非有限值 & 0 \checkmark \\
B & image\_\allowbreak{}deconvolution/image\_\allowbreak{}data 全有限 & 0 个非有限值 & 0 \checkmark \\
B & image\_\allowbreak{}deconvolution/image\_\allowbreak{}recon 全有限 & 0 个非有限值 & 0 \checkmark \\
B & image\_\allowbreak{}deconvolution/image\_\allowbreak{}abs\_\allowbreak{}error 全有限 & 0 个非有限值 & 0 \checkmark \\
B & image\_\allowbreak{}deconvolution/image\_\allowbreak{}residual 全有限 & 0 个非有限值 & 0 \checkmark \\
B & image\_\allowbreak{}deconvolution/history\_\allowbreak{}k 全有限 & 0 个非有限值 & 0 \checkmark \\
B & image\_\allowbreak{}deconvolution/error\_\allowbreak{}history 全有限 & 0 个非有限值 & 0 \checkmark \\
B & image\_\allowbreak{}deconvolution/residual\_\allowbreak{}history 全有限 & 0 个非有限值 & 0 \checkmark \\
B & image\_\allowbreak{}deconvolution/log10\_\allowbreak{}sigma 全有限 & 0 个非有限值 & 0 \checkmark \\
B & image\_\allowbreak{}deconvolution/log10\_\allowbreak{}lam\_\allowbreak{}sweep 全有限 & 0 个非有限值 & 0 \checkmark \\
B & image\_\allowbreak{}deconvolution/error\_\allowbreak{}lam\_\allowbreak{}sweep 全有限 & 0 个非有限值 & 0 \checkmark \\
B & bias\_\allowbreak{}variance/log10\_\allowbreak{}lam 全有限 & 0 个非有限值 & 0 \checkmark \\
B & bias\_\allowbreak{}variance/bias2\_\allowbreak{}analytic 全有限 & 0 个非有限值 & 0 \checkmark \\
B & bias\_\allowbreak{}variance/variance\_\allowbreak{}analytic 全有限 & 0 个非有限值 & 0 \checkmark \\
B & bias\_\allowbreak{}variance/mse\_\allowbreak{}analytic 全有限 & 0 个非有限值 & 0 \checkmark \\
B & bias\_\allowbreak{}variance/bias2\_\allowbreak{}mc 全有限 & 0 个非有限值 & 0 \checkmark \\
B & bias\_\allowbreak{}variance/variance\_\allowbreak{}mc 全有限 & 0 个非有限值 & 0 \checkmark \\
B & bias\_\allowbreak{}variance/mse\_\allowbreak{}mc 全有限 & 0 个非有限值 & 0 \checkmark \\
B & bias\_\allowbreak{}variance/mse\_\allowbreak{}mc\_\allowbreak{}std 全有限 & 0 个非有限值 & 0 \checkmark \\
B & bias\_\allowbreak{}variance/error\_\allowbreak{}rel\_\allowbreak{}analytic 全有限 & 0 个非有限值 & 0 \checkmark \\
B & bias\_\allowbreak{}variance/error\_\allowbreak{}rel\_\allowbreak{}mc 全有限 & 0 个非有限值 & 0 \checkmark \\
B & bias\_\allowbreak{}variance/dof 全有限 & 0 个非有限值 & 0 \checkmark \\
B & bias\_\allowbreak{}variance/local\_\allowbreak{}slope 全有限 & 0 个非有限值 & 0 \checkmark \\
B & bias\_\allowbreak{}variance/local\_\allowbreak{}slope\_\allowbreak{}variance 全有限 & 0 个非有限值 & 0 \checkmark \\
A & 四个可视化的数组数 / 构建耗时 / JSON 合法性 & pack\_frame 要求 meta 只有 number/\allowbreak{}bool/\allowbreak{}null、\allowbreak{}数组全有限、\allowbreak{}字符串走 strings & lcurve\_tradeoff: 17 数组 /\allowbreak{} 44 meta /\allowbreak{} 0.01s；filter\_shapes: 16 数组 /\allowbreak{} 30 meta /\allowbreak{} 0.04s；image\_deconvolution: 11 数组 /\allowbreak{} 27 meta /\allowbreak{} 0.23s；bias\_variance: 13 数组 /\allowbreak{} 27 meta /\allowbreak{} 0.04s(\allowbreak{}最慢 0.23s) \checkmark \\
A & 参数角点 + t 端点的帧数与最慢单帧耗时 / 非有限 meta 数 & 奇异状态必须变成可显示的值(\allowbreak{}None + strings 说明)，不能是 NaN/\allowbreak{}Inf & 67 帧全部通过 pack\_frame；最慢单帧 0.40s；非有限 meta 0 个 \checkmark \\
A & 帧 meta 与内核重算的最大相对差 & 帧的每个读数都由内核函数直接给出，帧层不做第二套计算 & max rel diff =\allowbreak{} 0.000e+00；log10\_lam\_oracle: -1.8；log10\_lam\_lcurve: -3.13333；log10\_lam\_gcv: -2.33333；error\_oracle: 0.00772735 \checkmark \\
A & 同自由度(dof = k\_\allowbreak{}TSVD)下的三方误差与赢家 & 形状不同不等于精度不同：必须在同一条自由度轴上比误差 & matched\_dof =\allowbreak{} 8；tsvd 0.00350，tikhonov 0.00283，landweber 0.09240 \checkmark \\
A & 误差-自由度曲线的形状：TSVD 是离散点、Tikhonov 是连续曲线 & TSVD 的自由度只能取整数 k；Tikhonov 的 lambda 连续 =\allowbreak{}> 自由度连续 & TSVD 曲线 20 点(\allowbreak{}自由度 1..20)；Tikhonov 曲线 61 点(\allowbreak{}自由度 0.01..19.99) \checkmark \\
A & 发散帧与正常帧的掩码零点数 / 因子越界幅度 & omega\_rel > 2 时 |1 - omega sigma\textasciicircum{}2| > 1，滤波因子离开 [0,\allowbreak{}1]，帧里被夹住并标记 & omega\_rel=\allowbreak{}2.5: max|f| =\allowbreak{} 1.918e+05，掩码 0 的个数 1；omega\_rel=\allowbreak{}1.0: 掩码 0 的个数 0 \checkmark \\
A & 发散状态在 strings 里有说明(不是只留一个数) & 占位值三件套：地板/\allowbreak{}夹住的值 + 掩码 + 文案 & omega\_note 提到越界与掩码：True \checkmark \\
A & 两种角度下的零空间维数 / 无噪时的误差地板 & A 是 m x n 时 rank <=\allowbreak{} m；m < n 必有零空间，零空间分量无法从数据恢复 & n\_angles=\allowbreak{}8: null =\allowbreak{} 20，误差 0.1577；n\_angles=\allowbreak{}16: null =\allowbreak{} 0，误差 4.28e-08 \checkmark \\
A & 秩亏与良态两种状态在 strings 里都能读到 & 用户看到的是图，图注必须能解释误差为什么不再下降 & null\_space\_note 提到零空间：True \checkmark \\
A & 帧内解析 / 蒙特卡洛总误差曲线的最大相对差 & 64 个噪声样本的经验均值误差 \textasciitilde{} sigma/\allowbreak{}sqrt(\allowbreak{}64) =\allowbreak{}> 相对差 10\% 量级 & max rel diff =\allowbreak{} 0.099；n\_trials =\allowbreak{} 64 \checkmark \\
A & 大方差端的对数-对数斜率 meta(理论 -4) & lambda >> sigma\_1 时 Var \textasciitilde{} lambda\textasciicircum{}(\allowbreak{}-4) & variance\_slope\_high =\allowbreak{} -3.9993 \checkmark \\
A & 两侧失败方式的读数(lambda 最小 / 最大处的相对误差) & 左端噪声接管、\allowbreak{}右端解被压成 0；两者都应当明显大于最小值 & error\_at\_lam\_min =\allowbreak{} 1.529，error\_at\_lam\_max =\allowbreak{} 1.269，最小误差 0.0199 \checkmark \\
A & 对数轴曲线的最小值(都应当 >= 地板) & 对数轴上 0 画不出来，曲线抬到 1e-17 并说明 & log10\_residual 7.012e-04，error\_curve 7.727e-03，gcv 1.307e-09，curvature 1.000e-17 \checkmark \\
A & Picard 系数落到噪声之下的掩码(0/1 且有 0 也有 1) & 噪声接管之后 |u\_i\textasciicircum{}T b| < ||e||，这一段不能用来判断 Picard 条件 & 掩码和 =\allowbreak{} 15/\allowbreak{}20(\allowbreak{}meta 记 15) \checkmark \\
A & null(无有限值)在 strings 里有专门说明 & meta 里的 null 必须能解释成没测到，而不是 0 & null\_note 存在且含 null：True \checkmark \\
B & lcurve\_\allowbreak{}tradeoff 提供 requiredAux=error & $\le$ 2 & 1 \checkmark \\
B & filter\_\allowbreak{}shapes 提供 requiredAux=error & $\le$ 4 & 1 \checkmark \\
B & image\_\allowbreak{}deconvolution 提供 requiredAux=error & $\le$ 3 & 1 \checkmark \\
B & image\_\allowbreak{}deconvolution 提供 requiredAux=residual & $\le$ 2 & 1 \checkmark \\
B & bias\_\allowbreak{}variance 提供 requiredAux=error & $\le$ 2 & 1 \checkmark \\
B & lcurve\_\allowbreak{}tradeoff 的 theory 条数 & $\le$ 3 & 1 \checkmark \\
B & lcurve\_\allowbreak{}tradeoff 的 pitfalls 长度 & $\le$ 359 & 60 \checkmark \\
B & lcurve\_\allowbreak{}tradeoff 的 pitfalls 含数字 & $\le$ 32 & 1 \checkmark \\
B & lcurve\_\allowbreak{}tradeoff/log10\_\allowbreak{}lam 的 criticalNote 非空 & $\le$ 120 & 1 \checkmark \\
B & filter\_\allowbreak{}shapes 的 theory 条数 & $\le$ 3 & 1 \checkmark \\
B & filter\_\allowbreak{}shapes 的 pitfalls 长度 & $\le$ 307 & 60 \checkmark \\
B & filter\_\allowbreak{}shapes 的 pitfalls 含数字 & $\le$ 10 & 1 \checkmark \\
B & filter\_\allowbreak{}shapes/log10\_\allowbreak{}lam 的 criticalNote 非空 & $\le$ 87 & 1 \checkmark \\
B & filter\_\allowbreak{}shapes/omega\_\allowbreak{}rel 的 criticalNote 非空 & $\le$ 157 & 1 \checkmark \\
B & image\_\allowbreak{}deconvolution 的 theory 条数 & $\le$ 3 & 1 \checkmark \\
B & image\_\allowbreak{}deconvolution 的 pitfalls 长度 & $\le$ 271 & 60 \checkmark \\
B & image\_\allowbreak{}deconvolution 的 pitfalls 含数字 & $\le$ 14 & 1 \checkmark \\
B & image\_\allowbreak{}deconvolution/blur\_\allowbreak{}sigma 的 criticalNote 非空 & $\le$ 95 & 1 \checkmark \\
B & bias\_\allowbreak{}variance 的 theory 条数 & $\le$ 3 & 1 \checkmark \\
B & bias\_\allowbreak{}variance 的 pitfalls 长度 & $\le$ 264 & 60 \checkmark \\
B & bias\_\allowbreak{}variance 的 pitfalls 含数字 & $\le$ 16 & 1 \checkmark \\
B & bias\_\allowbreak{}variance/log10\_\allowbreak{}lam 的 criticalNote 非空 & $\le$ 68 & 1 \checkmark \\
B & lcurve\_\allowbreak{}tradeoff/log10\_\allowbreak{}lam critical >= min & $\le$ 0 & 0 \checkmark \\
B & lcurve\_\allowbreak{}tradeoff/log10\_\allowbreak{}lam critical <= max & $\le$ 0 & 0 \checkmark \\
B & filter\_\allowbreak{}shapes/log10\_\allowbreak{}lam critical >= min & $\le$ 0 & 0 \checkmark \\
B & filter\_\allowbreak{}shapes/log10\_\allowbreak{}lam critical <= max & $\le$ 0 & 0 \checkmark \\
B & filter\_\allowbreak{}shapes/omega\_\allowbreak{}rel critical >= min & $\le$ 0 & 0 \checkmark \\
B & filter\_\allowbreak{}shapes/omega\_\allowbreak{}rel critical <= max & $\le$ 0 & 0 \checkmark \\
B & filter\_\allowbreak{}shapes/omega\_\allowbreak{}rel critical >= min & $\le$ 0 & 0 \checkmark \\
B & filter\_\allowbreak{}shapes/omega\_\allowbreak{}rel critical <= max & $\le$ 0 & 0 \checkmark \\
B & image\_\allowbreak{}deconvolution/blur\_\allowbreak{}sigma critical >= min & $\le$ 0 & 0 \checkmark \\
B & image\_\allowbreak{}deconvolution/blur\_\allowbreak{}sigma critical <= max & $\le$ 0 & 0 \checkmark \\
B & bias\_\allowbreak{}variance/log10\_\allowbreak{}lam critical >= min & $\le$ 0 & 0 \checkmark \\
B & bias\_\allowbreak{}variance/log10\_\allowbreak{}lam critical <= max & $\le$ 0 & 0 \checkmark \\
C & 最贵参数组合下的单帧耗时 & 交互式可视化要求单帧 <=\allowbreak{} 5s(\allowbreak{}预算)，帧层不得有隐藏的 O(\allowbreak{}n\textasciicircum{}3) 重算 & lcurve\_tradeoff 0.01s；filter\_shapes 0.04s；image\_deconvolution 0.50s；bias\_variance 0.31s(\allowbreak{}最慢 0.50s) \checkmark \\
A & 曲率拐点下标(本模块 vs 1.3 lcurve\_\allowbreak{}corner) & 同一离散曲率公式(\allowbreak{}相邻三点中心差分)=\allowbreak{}> 下标必须逐位相同 & noise=\allowbreak{}1e-2: 37 vs 37；noise=\allowbreak{}1e-3: 29 vs 29；noise=\allowbreak{}1e-4: 22 vs 22；noise=\allowbreak{}1e-5: 14 vs 14；noise=\allowbreak{}1e-6: 4 vs 4 \checkmark \\
A & 三个拐点判据的 lambda 极差 / 下标跨度 & 曲率、\allowbreak{}两段直线交点、\allowbreak{}离散二阶差分是三个不同启发式，没有理由给出同一个数 & kappa=\allowbreak{}1e3: lambda 极差 681 倍、\allowbreak{}下标跨度 34(\allowbreak{}曲率 idx 5，两段直线 idx 39，二阶差分 idx 26)；kappa=\allowbreak{}1e4: lambda 极差 126 倍、\allowbreak{}下标跨度 21(\allowbreak{}曲率 idx 20，两段直线 idx 41，二阶差分 idx 32)；kappa=\allowbreak{}1e5: lambda 极差 33 倍、\allowbreak{}下标跨度 13(\allowbreak{}曲率 idx 24，两段直线 idx 31，二阶差分 idx 37)；kappa=\allowbreak{}1e6: lambda 极差 40 倍、\allowbreak{}下标跨度 12(\allowbreak{}曲率 idx 29，两段直线 idx 41，二阶差分 idx 40)；kappa=\allowbreak{}1e7: lambda 极差 63 倍、\allowbreak{}下标跨度 12(\allowbreak{}曲率 idx 33，两段直线 idx 45，二阶差分 idx 42) \checkmark \\
A & 无约束 vs 有约束两段直线拟合：一致组数 / 最大偏离倍数(18 组配置) & 无约束时段越短拟合残差越小 =\allowbreak{}> 切点无规律漂移；约束把切点限制在 [25\%,\allowbreak{} 75\%] 且交点必须落在包围盒内 & 一致 7/\allowbreak{}18 组；其余偏离最大 125.9 倍(\allowbreak{}中位 1.39) \checkmark \\
A & 有约束拐点的 lambda 落在 L 曲线拐弯区间内 & 约束版的交点应当落在曲率极大点附近(\allowbreak{}同一段弯)，偏离量就是这条判据的不确定度 & 有约束 lambda =\allowbreak{} 2.929e-02，曲率判据 lambda =\allowbreak{} 7.356e-04，比值 39.81 \checkmark \\
A & GCV 曲线与 1.3 cross\_\allowbreak{}validation\_\allowbreak{}curve 的最大相对差 & G(\allowbreak{}lambda) =\allowbreak{} ||A x\_lambda - b||\textasciicircum{}2 /\allowbreak{} (\allowbreak{}m - sum f\_i)\textasciicircum{}2，两条实现同一个公式 & max rel diff =\allowbreak{} 3.502e-12；最小值位置 0.00522335 vs 0.00522335 \checkmark \\
A & GCV 不被噪声水平骗：lambda\_\allowbreak{}GCV 随噪声下降而下降 & 噪声越小，能安全保留的模态越多 =\allowbreak{}> 最优 lambda 越小 & lambda\_GCV =\allowbreak{} [0.016 0.005 0.002 0.001 0.   ](\allowbreak{}噪声 1e-2 -> 1e-6) \checkmark \\
A & lambda\_\allowbreak{}GCV 与 GCV 曲线最小点的相对差 & lambda\_GCV =\allowbreak{} argmin G(\allowbreak{}lambda)(\allowbreak{}定义) & |lambda\_GCV - lambda\_argmin| /\allowbreak{} lambda\_GCV =\allowbreak{} 0.000e+00 \checkmark \\
A & 偏差原则处的残差 / 噪声水平 & ||A x\_lambda - b|| =\allowbreak{} tau ||e||(\allowbreak{}tau =\allowbreak{} 1)，交点用对数插值 & lambda\_disc =\allowbreak{} 1.3396e-02，比值 =\allowbreak{} 0.985457 \checkmark \\
A & oracle 位置与误差曲线最小点 & lambda\_oracle =\allowbreak{} argmin ||x\_lambda - x*||(\allowbreak{}定义) & |lambda\_oracle - lambda\_argmin| /\allowbreak{} lambda\_oracle =\allowbreak{} 0.000e+00；最小误差 0.0077 \checkmark \\
A & 四个判据的误差放大倍数(oracle / 偏差原则 / GCV / L 曲线 / 不正则化) & 判据越依赖残差小这种间接信息，越偏向欠正则化的一端 & 1.000(\allowbreak{}oracle)| 1.159(\allowbreak{}disc)| 3.791(\allowbreak{}GCV)| 21.478(\allowbreak{}L)| 12863(\allowbreak{}无) \checkmark \\
A & 网格两端(lambda 最小 / 最大)的相对误差与解范数比 & 左端：无正则化的噪声放大；右端：x\_lambda -> 0 =\allowbreak{}> 误差 -> ||x*|| & 左端 [1.241e-01 1.115e+00 1.040e+01 9.940e+01 9.673e+02]；右端 [0.993 0.993 0.992 0.992 0.992]；||x(\allowbreak{}lambda\_max)||/\allowbreak{}||x(\allowbreak{}lambda\_min)|| =\allowbreak{} [7.50e-03 5.35e-03 8.05e-04 8.78e-05 9.28e-06] \checkmark \\
A & 曲线坐标与逐点独立求解的最大相对差(||x|| 与相对残差) & L 曲线的每个点就是该 lambda 处 Tikhonov 解的范数与残差 & ||x|| 相对差 0.000e+00，相对残差 0.000e+00 \checkmark \\
A & current\_\allowbreak{}solution 与 tikhonov 的一致性(拖动 lambda 的那条路) & 交互拖动 lambda 时前端拿到的是同一个解 & max|dx| =\allowbreak{} 0.000e+00 \checkmark \\
B & 三个拐点下标都在网格内部(kappa=1e3) & $\le$ 0 & 0 \checkmark \\
B & 三个拐点下标都在网格内部(kappa=1e5) & $\le$ 0 & 0 \checkmark \\
B & 三个拐点下标都在网格内部(kappa=1e7) & $\le$ 0 & 0 \checkmark \\
B & 有效判据读数全有限 & 0 个非有限值 & 0 \checkmark \\
B & 无噪时偏差原则没有交点 -> NaN(不是 0) & 1 & 1 \checkmark \\
B & ||x\_\allowbreak{}lambda|| 随 lambda 不增 & nonincreasing & 违反 0 处 \checkmark \\
B & 相对残差随 lambda 不减 & nondecreasing & 违反 0 处 \checkmark \\
B & 离散曲率全有限 & 0 个非有限值 & 0 \checkmark \\
B & 网格左端 = sigma\_\allowbreak{}n/10 & $1.000\times 10^{-7}$ & $1.000\times 10^{-7}$ \checkmark \\
B & 网格右端 = 10 sigma\_\allowbreak{}1 & 10 & 10 \checkmark \\
B & 网格单调递增 & increasing & 违反 0 处 \checkmark \\
C & L 曲线代价降到 10 倍以下的噪声水平(二分) & 最优 lambda 随噪声下降，而 L 曲线拐点滞后 =\allowbreak{}> 噪声越小拐点越差 & 阈值 =\allowbreak{} 4.1727e-03(\allowbreak{}对数 2.38 个数量级)；噪声 1e-2..1e-6 时的代价 [  6.5  21.5  30.  199.1 211.6] \checkmark \\
C & lambda\_\allowbreak{}GCV 对噪声水平的幂律指数 & 几何谱 + GCV：lambda\_GCV \textasciitilde{} noise\textasciicircum{}(\allowbreak{}1/\allowbreak{}2)(\allowbreak{}实测指数略低于 1/\allowbreak{}2) & p =\allowbreak{} 0.4667(\allowbreak{}R\textasciicircum{}2 =\allowbreak{} 0.997557)；lambda\_GCV =\allowbreak{} [0.02 0.   0.   0.   0.  ] \checkmark \\
C & lambda\_\allowbreak{}oracle 对噪声水平的幂律指数(对照) & oracle 的 lambda 也随噪声下降；GCV 的指数与它的差就是判据的滞后量 & p =\allowbreak{} 0.4800(\allowbreak{}GCV 比它小 0.013) \checkmark \\
C & 网格左端误差对 kappa 的幂律指数 & 不做正则化时误差 \textasciitilde{} ||e||/\allowbreak{}(\allowbreak{}sigma\_n ||x*||) \textasciitilde{} kappa & p =\allowbreak{} 0.9733(\allowbreak{}R\textasciicircum{}2 =\allowbreak{} 0.999950)；误差 [1.241e-01 1.115e+00 1.040e+01 9.940e+01 9.673e+02] \checkmark \\
A & 滤波因子与闭式的最大偏差 / 半功率点 & f\_i =\allowbreak{} s\_i\textasciicircum{}2/\allowbreak{}(\allowbreak{}s\_i\textasciicircum{}2+lambda\textasciicircum{}2)；s\_i =\allowbreak{} lambda 处 f =\allowbreak{} 1/\allowbreak{}2 & max|df| =\allowbreak{} 0.000e+00；f(\allowbreak{}s=\allowbreak{}lambda) =\allowbreak{} 0.500000000000 \checkmark \\
A & lambda=0 与最小二乘解的一致性 / lambda 极大时的压缩比 & lambda=\allowbreak{}0 时 Tikhonov =\allowbreak{} 最小二乘；lambda >> sigma\_1 时 ||x|| \textasciitilde{} ||x\_ls|| sigma\_1\textasciicircum{}2/\allowbreak{}lambda\textasciicircum{}2 & rel diff =\allowbreak{} 2.819e-11；||x\_huge||/\allowbreak{}||x\_ls|| =\allowbreak{} 8.796e-15 \checkmark \\
A & 谱公式与 LU 正规方程解的最大相对差 & <=\allowbreak{} 1e-10(\allowbreak{}lambda 不太小时两条路线同一个解；kappa=\allowbreak{}1e6 下实测 2e-12) & lambda=\allowbreak{}0.01: 1.97e-12；lambda=\allowbreak{}0.1: 1.91e-14；lambda=\allowbreak{}1: 7.36e-16；lambda=\allowbreak{}0.01: 1.97e-12 \checkmark \\
A & 正规方程路线在 lambda -> 0 时的退化(kappa 平方的代价) & 正规方程的条件数是 kappa(\allowbreak{}A)\textasciicircum{}2；lambda 小到正则项压不住时 LU 解开始丢位 & lambda=\allowbreak{}1e-2: 1.97e-12；lambda=\allowbreak{}1e-6: 3.19e-06(\allowbreak{}劣化 1614925 倍；kappa(\allowbreak{}A)\textasciicircum{}2 =\allowbreak{} 1.000e+12) \checkmark \\
A & TSVD 与 1.3 truncated\_\allowbreak{}solution 的最大逐点差 / 误差曲线 & k =\allowbreak{} n 时退化为最小二乘解；误差在 k* 处最小、\allowbreak{}之后被噪声抬高 & max|dx| =\allowbreak{} 5.551e-17；k* =\allowbreak{} 7(\allowbreak{}err 0.0163)，k =\allowbreak{} n 时 err =\allowbreak{} 100.1523(\allowbreak{}放大 6141.3 倍) \checkmark \\
A & 谱形式与矩阵迭代的最大相对差 / 发散线另一侧的表现 & 两者同解(\allowbreak{}<=\allowbreak{} 1e-9)；omega >=\allowbreak{} 2/\allowbreak{}sigma\_1\textasciicircum{}2 时迭代发散、\allowbreak{}滤波因子离开 [0,\allowbreak{}1] & omega\_rel=\allowbreak{}1,\allowbreak{}k=\allowbreak{}30: 6.63e-14(\allowbreak{}stop=\allowbreak{}n\_iter)；omega\_rel=\allowbreak{}0.5,\allowbreak{}k=\allowbreak{}60: 2.61e-13(\allowbreak{}stop=\allowbreak{}n\_iter)；omega\_rel=\allowbreak{}1.5,\allowbreak{}k=\allowbreak{}12: 3.18e-14(\allowbreak{}stop=\allowbreak{}n\_iter)；omega\_rel=\allowbreak{}2.2: stop=\allowbreak{}diverged，max f =\allowbreak{} 6.859e+15 \checkmark \\
A & 两点 TV 闭式解的最大逐点偏差 & x =\allowbreak{} (\allowbreak{}m - soft(\allowbreak{}d,\allowbreak{} lambda),\allowbreak{} m + soft(\allowbreak{}d,\allowbreak{} lambda))，m =\allowbreak{} mean(\allowbreak{}b)，d =\allowbreak{} (\allowbreak{}b\_2-b\_1)/\allowbreak{}2 & lambda=\allowbreak{}0: max|dx|=\allowbreak{}0.00e+00；lambda=\allowbreak{}0.05: max|dx|=\allowbreak{}1.11e-16；lambda=\allowbreak{}0.3: max|dx|=\allowbreak{}1.11e-16；lambda=\allowbreak{}0.75: max|dx|=\allowbreak{}5.55e-17；lambda=\allowbreak{}2: max|dx|=\allowbreak{}1.11e-16 \checkmark \\
A & lambda >= |d| 时 TV 解是否退化成常数 & TV 罚项在差值上的次梯度包含 0 时，两点被合并(\allowbreak{}台阶先验) & lambda=\allowbreak{}1.0: x =\allowbreak{} (\allowbreak{}0.250000000,\allowbreak{} 0.250000000)，差 =\allowbreak{} 1.11e-16 \checkmark \\
A & lambda=0 时 TV 探针是否退化为恒等(正对照，kappa=1e2) & lambda=\allowbreak{}0 时 TV 项不起作用，x =\allowbreak{} A\textasciicircum{}+ b，探针读数应为 1 & lambda=\allowbreak{}0 的探针读数最小偏差 =\allowbreak{} 5.45e-14(\allowbreak{}2000 步) \checkmark \\
A & TV 探针的依赖性 / 非对角泄漏 & TV 非线性且非对角：换探针读数会变，泄漏不为 0(\allowbreak{}它不是滤波因子) & max|f\_mixed - f\_single| =\allowbreak{} 0.959；max leakage =\allowbreak{} 0.144 \checkmark \\
A & 蒙特卡洛总误差落在解析值的 4 倍标准误内 & 64 个样本的经验均值误差 \textasciitilde{} sigma/\allowbreak{}sqrt(\allowbreak{}64)；逐点比较要用标准误刻度 & 越界点 0/\allowbreak{}41；最大相对差 mse 0.099，variance 0.118 \checkmark \\
A & 偏差\textasciicircum{}2 + 方差 = 总误差(恒等式) & 两项按定义相加就是 E||x\_lambda - x*||\textasciicircum{}2 & max|bias2+variance-mse| =\allowbreak{} 0.000e+00 \checkmark \\
A & lambda 最小时实测偏差\textasciicircum{}2 与解析偏差\textasciicircum{}2 的绝对差(以 mse\_\allowbreak{}min 为单位) & 经验均值 E\_hat[x] 的抽样误差是 sigma/\allowbreak{}sqrt(\allowbreak{}N)，它进入偏差\textasciicircum{}2 后不可忽略 & bias2\_analytic =\allowbreak{} 6.407e-09，bias2\_mc =\allowbreak{} 4.354e-02，绝对差 =\allowbreak{} 4.354e-02 =\allowbreak{} 106.5 x mse\_min \checkmark \\
A & 噪声对齐最小奇异方向时的相对误差与解析值 & ||A\textasciicircum{}+ e|| =\allowbreak{} ||e||/\allowbreak{}sigma\_n(\allowbreak{}e 沿 u\_n)，相对误差 =\allowbreak{} ||e||/\allowbreak{}(\allowbreak{}sigma\_n ||x*||) & 实测 900.363172 /\allowbreak{} 解析 900.363172(\allowbreak{}相对差 7.92e-12) \checkmark \\
A & kappa=1e6、噪声 1e-3 时的放大倍数(最坏 / 均方根) & 最坏 ||e||sqrt(\allowbreak{}sum 1/\allowbreak{}s\_i\textasciicircum{}2)/\allowbreak{}||x*||，均方根除以 sqrt(\allowbreak{}m) & 最坏 1028.4 倍，均方根 181.8 倍(\allowbreak{}两者相差 5.7 倍) \checkmark \\
A & 无噪 Picard 斜率 / 比值斜率 / 相关系数 & alpha\_i =\allowbreak{} sigma\_i\textasciicircum{}decay(\allowbreak{}decay=\allowbreak{}2)=\allowbreak{}> 斜率 =\allowbreak{} 2，比值斜率 =\allowbreak{} 1，log-log 完全共线 & slope =\allowbreak{} 2.000000，ratio slope =\allowbreak{} 1.000000，corr =\allowbreak{} 1.000000 \checkmark \\
A & 有噪时整段斜率 vs 信号段斜率 / 被噪声接管的模态数 & 噪声接管后 |u\_i\textasciicircum{}T b| 不再随 sigma\_i 下降，整段斜率必然变平；信号段仍满足理论 & 整段 0.5641，信号段 1.9972(\allowbreak{}5 个模态)，低于噪声的模态 15(\allowbreak{}按 sigma\_res 计 10) \checkmark \\
A & Picard 不成立时：oracle 位置与正则化的收益 & |u\_i\textasciicircum{}T b|/\allowbreak{}sigma\_i 随 i 增长时，最小二乘解在噪声下必然崩，正则化只能截掉一部分 & oracle lambda =\allowbreak{} 1.000e-07(\allowbreak{}网格左端 1.000e-07)，误差 0.116，正则化收益 =\allowbreak{} 1.000 倍 \checkmark \\
A & Picard 成立时的对照：最优 lambda 处误差与不正则化的比 & 同一算子、\allowbreak{}同一噪声，Picard 成立时正则化把误差从 1e4 量级压回 1e-2 量级 & 不正则化误差 99.399，oracle 处 0.0077，比值 12863 倍 \checkmark \\
A & 二维 Tikhonov 解的正规方程残差(相对) & (\allowbreak{}B\textasciicircum{}T B) X (\allowbreak{}B\textasciicircum{}T B) + lambda\textasciicircum{}2 X =\allowbreak{} B\textasciicircum{}T Y B 是二维 Tikhonov 的定义方程 & max|残差|/\allowbreak{}max|B\textasciicircum{}T Y B| =\allowbreak{} 5.130e-15 \checkmark \\
A & 二维条件数 = 一维条件数的平方 & 可分离算子的奇异值是 s\_i s\_j，因此 kappa\_2D =\allowbreak{} kappa\_1D\textasciicircum{}2 & kappa\_1D =\allowbreak{} 5911 -> kappa\_2D =\allowbreak{} 3.494e+07(\allowbreak{}s\_1 s\_1 =\allowbreak{} 1，s\_n s\_n =\allowbreak{} 2.862e-08) \checkmark \\
A & 停滞 SVD 与体检后 SVD 的相对误差(对 LAPACK) & 单边 Jacobi 在重奇异值矩阵上可能停滞；factorize 用 converged + 重建 + 正交性三道体检回退 & jacobi: converged=\allowbreak{}False，相对误差 2221.2 倍；factorize -> kogbetliantz(\allowbreak{}fallback=\allowbreak{}1)，相对误差 1.88e-13 \checkmark \\
A & CT 投影矩阵的行和 / 伴随恒等式 / 两种角度下的秩 & 射线驱动的线积分归一化；A 的秩由角度数决定，角度不足时零空间非平凡 & 每行行和与该角度弦长的最大偏差 5.329e-15(\allowbreak{}8 个角度 x 10 条射线)；|<Ax,\allowbreak{}y>-<x,\allowbreak{}A\textasciicircum{}T y>| =\allowbreak{} 5.773e-15；8 角度秩 80/\allowbreak{}100，14 角度秩 100/\allowbreak{}100，kappa 从 1.97e+05 降到 1.78e+01 \checkmark \\
B & ||x\_\allowbreak{}lambda|| 随 lambda 不增 & nonincreasing & 违反 0 处 \checkmark \\
B & 相对残差随 lambda 不减 & nondecreasing & 违反 0 处 \checkmark \\
B & 有效自由度随 lambda 不增 & nonincreasing & 违反 0 处 \checkmark \\
B & TSVD 残差随 k 不增 & nonincreasing & 违反 0 处 \checkmark \\
B & TSVD 滤波因子只有 0/1 & 0 & 0 \checkmark \\
B & Landweber 上界(omega\_\allowbreak{}rel=0.1) & $\le$ 1 & 1 \checkmark \\
B & Landweber 下界(omega\_\allowbreak{}rel=0.1) & $\le$ 0 & 0 \checkmark \\
B & Landweber 上界(omega\_\allowbreak{}rel=0.5) & $\le$ 1 & 1 \checkmark \\
B & Landweber 下界(omega\_\allowbreak{}rel=0.5) & $\le$ 0 & 0 \checkmark \\
B & Landweber 上界(omega\_\allowbreak{}rel=1.0) & $\le$ 1 & 1 \checkmark \\
B & Landweber 下界(omega\_\allowbreak{}rel=1.0) & $\le$ 0 & 0 \checkmark \\
B & Landweber 上界(omega\_\allowbreak{}rel=1.9) & $\le$ 1 & 1 \checkmark \\
B & Landweber 下界(omega\_\allowbreak{}rel=1.9) & $\le$ 0 & 0 \checkmark \\
B & 自由度上界 & $\le$ 20 & 10 \checkmark \\
B & 最优步之前的误差不增 & nonincreasing & 违反 0 处 \checkmark \\
B & 最优步处误差小于末步(半收敛) & $\le$ 0.035007 & 0.035007 \checkmark \\
B & 误差历史全有限 & 0 个非有限值 & 0 \checkmark \\
B & TV 目标函数不超过候选点 & $\le$ 0.0615792 & 0.0337634 \checkmark \\
B & 幻影只有四档灰度(含外围 1.0) & 0 & 0 \checkmark \\
B & 模糊矩阵行和为 1 & 0 & 0 \checkmark \\
B & 模糊矩阵对称 & 0 & 0 \checkmark \\
C & 矩阵迭代实测的发散阈值(400 步预算) & 2 & 2.07152 \checkmark \\
C & 谱判据给出的发散阈值(= 2/sigma\_\allowbreak{}1\textasciicircum{}2，精确) & 2 & 2 \checkmark \\
C & 方差尾部的幂律指数(lambda 从 1e5 降到 10，25 个点) & lambda >> sigma\_1 时 f\_i \textasciitilde{} s\_i\textasciicircum{}2/\allowbreak{}lambda\textasciicircum{}2 =\allowbreak{}> Var \textasciitilde{} lambda\textasciicircum{}(\allowbreak{}-4)，拟合区间 lambda in [10 sigma\_1,\allowbreak{} 1e5] & p =\allowbreak{} -3.9994(\allowbreak{}窗口 1e+01..1e+05，R\textasciicircum{}2 =\allowbreak{} 1.000000) \checkmark \\
C & 最坏放大倍数对 kappa 的幂律指数 & 最坏放大 \textasciitilde{} ||e|| sqrt(\allowbreak{}sum 1/\allowbreak{}s\_i\textasciicircum{}2)/\allowbreak{}||x*||；几何谱下 sum 1/\allowbreak{}s\_i\textasciicircum{}2 由 1/\allowbreak{}sigma\_n\textasciicircum{}2 主导 =\allowbreak{}> 指数 1 & p =\allowbreak{} 0.9780(\allowbreak{}R\textasciicircum{}2 =\allowbreak{} 0.999942)；放大倍数 [1.3e-01 1.1e+00 1.1e+01 1.0e+02 1.0e+03] \checkmark \\
C & 放大倍数与主项 1/sigma\_\allowbreak{}n 的比(多留的模态贡献多少) & 最坏放大 =\allowbreak{} ||e||/\allowbreak{}sigma\_n * sqrt(\allowbreak{}1 + sigma\_n\textasciicircum{}2 sum\_\{i<n\} 1/\allowbreak{}s\_i\textasciicircum{}2)，第二项随 kappa 递减 & 比值 [1.6134 1.3912 1.2693 1.1932 1.1423](\allowbreak{}kappa=\allowbreak{}1e2 高出 61\%，kappa=\allowbreak{}1e6 只高出 14\%) \checkmark \\
\end{longtable}


\section{它会怎样骗你（本章版）}

\begin{enumerate}
  \item \textbf{"L 曲线拐点就是最优 $\lambda$"}：$\kappa=10^6$、噪声 $10^{-3}$ 时拐点比
        oracle 小 $21.5$ 倍、误差放大 $21.5$ 倍；GCV 只放大 $3.79$ 倍。
  \item \textbf{"拐点是一个客观量"}：三个判据的 $\lambda$ 极差 $33\sim681$ 倍，下标跨度 $7\sim42$。
  \item \textbf{"无约束的两段直线拟合更客观"}：18 组里 11 组与带约束版不同，最大偏 $125.9$ 倍；
        我最初写在 docstring 里的"它总落在网格端点"实测不成立。
  \item \textbf{"GCV 曲线形状对就行"}：多乘一个 $\sigma$ 会让它错 $3.2\times10^4$ 倍，
        而且 $\lambda_{\mathrm{GCV}}$ 对噪声完全不敏感（斜率 $0.000$ 而不是 $1/2$）。
  \item \textbf{"帧里有掩码就算三件套齐了"}：掩码先夹后判会全是 $1$，
        发散段的占位值被当成真数据。
  \item \textbf{"用蒙特卡洛量偏差$^2$"}：小 $\lambda$ 端样本均值的抽样误差把偏差$^2$ 抬高
        $106.5$ 倍 ${\rm mse}_{\min}$；偏差要用无噪右端项算。
  \item \textbf{"正则化让结果更准"}：$\lambda$ 两端误差 $1.53$ 与 $1.27$，中间才有 $0.0199$；
        偏差与方差此消彼长，但"更准"只在中间一段成立。
  \item \textbf{"整段 Picard 斜率 $0.56$，所以不能正则化"}：信号段斜率 $1.997$；
        噪声接管之后整段斜率必然变平。
  \item \textbf{"Picard 不成立也能靠正则化救回来"}：此时最优 $\lambda$ 在网格左端，
        收益 $1.000$ 倍，即"不要正则化"。
  \item \textbf{"把 TV 画成第四条滤波曲线"}：$\lambda=0$ 时探针确实是恒等，
        但 $\lambda>0$ 时换探针读数变 $0.959$、泄漏 $0.144$。
  \item \textbf{"CT 是病态问题，误差降不下去是噪声的错"}：$8$ 条角度时零空间 $20$ 维，
        噪声拖到 $10^{-8}$ 误差仍 $0.158$；$14$ 条角度时它是良态的（$\kappa=17.8$）。
  \item \textbf{"残差小就是解准"}：Landweber 在 $k^\star$ 之后残差继续下降而误差上升；
        早停是正则化的一部分，不是停机技巧。
  \item \textbf{"同样自由度下三种方法形状不同，精度差不多"}：对齐自由度 $8$ 时
        TSVD $3.50\times10^{-3}$、Tikhonov $2.83\times10^{-3}$、Landweber $9.24\times10^{-2}$，
        相差 $33$ 倍，必须实测。
  \item \textbf{"阈值 $2/\sigma_1^2$ 抄书就行"}：矩阵迭代在 $400$ 步预算下实测阈值 $2.0715$，
        比理论高 $3.6\%$；两条阈值同时报才能说明"偏差是预算造成的"。
\end{enumerate}
